% TOP_PROBLEM: {"id": 1509, "title": "Siegel Zeros", "category_id": 1, "category": {"id": 1, "name": "number_theory", "display_name": "Number Theory", "description": "Properties of integers, prime numbers, Diophantine equations.", "slug": "number-theory", "order_index": 1, "created_at": "2026-07-31T15:26:25.670Z"}, "status": "open"}
% FIRST_POSTED: 2026-09-25
% ENTRY_KIND: statement_only
% !TeX program = lualatex
% Definition and sources only; this file is not an LLM solution attempt.
\documentclass[11pt]{article}
\usepackage[margin=25mm]{geometry}
\usepackage{fontspec}
\setmainfont{Latin Modern Roman}
\usepackage{amsmath,amssymb,mathtools}
\usepackage{unicode-math}
\setmathfont{Latin Modern Math}
\usepackage{xurl,hyperref}
\hypersetup{colorlinks=true,urlcolor=blue}
\setcounter{secnumdepth}{0}
\setlength{\parindent}{0pt}
\setlength{\parskip}{5pt}
\setlength{\emergencystretch}{3em}
\begin{document}
\section*{\texorpdfstring{Siegel Zeros}{Siegel Zeros}}\label{1509}
\subsection{Definitions and mathematical statement}
A primitive quadratic character $\chi$ has minimal modulus, its conductor $q$, and values in $\{0,\pm1\}$. Using its Dirichlet $L$-function, the uniform zero-free assertion is
\[
 \exists\delta>0\ \forall q\ge3\ \forall\chi\text{ primitive quadratic of conductor }q,
 \quad L(\beta,\chi)\ne0\quad\bigl(1-\delta/\log q<\beta<1\bigr).
\]
The same $\delta$ must work for all conductors, and only real $\beta$ are quantified here. A separate positive zero-free radius for each individual character would not establish the target.
\par\smallskip\textit{Formulation and concept references:} \href{https://arxiv.org/pdf/2602.03626v1}{[S1]}.

\subsection{Short English statement}
Can one uniformly exclude exceptionally close real zeros of quadratic Dirichlet $L$-functions near one?
\subsection{Sources}
\begin{itemize}
\item[S1]Rick F. Lu, Asif Zaman and Haonan Zhao, Numerical Computations Concerning Landau–Siegel Zeros; arXiv2602.03626v1 February3,2026.
\url{https://arxiv.org/pdf/2602.03626v1}.
\textit{Formulation source.}
\item[S2]On the Extended, Generalized, and Grand Riemann Hypotheses: A Unified Approach via General Properties of L-Functions, v18.
\url{https://www.preprints.org/manuscript/202506.0481/v18}.
\textit{Further reference; not a proof endorsement.}
\item[S3]Numerical Computations concerning Landau–Siegel Zeros.
\url{https://arxiv.org/html/2602.03626v1}.
\textit{Further reference; not a proof endorsement.}
\item[S4]Landau-Siegel zeros of Rankin-Selberg L-functions.
\url{https://arxiv.org/abs/2601.04189}.
\textit{Further reference; not a proof endorsement.}
\end{itemize}
\end{document}
